\quad O Insertion Sort é um algoritmo de ordenação baseado na comparação
e inserção dos elementos em suas posições corretas. Seu funcionamento é
semelhante ao ato de organizar cartas de baralho: os elementos são
analisados um por vez e cada novo elemento é inserido na posição adequada
da parte que já está ordenada.

\quad Inicialmente, considera-se que o primeiro elemento do vetor já está
ordenado. Em seguida, o algoritmo seleciona o próximo elemento, chamado de
chave, e compara seu valor com os elementos anteriores. Enquanto esses
elementos forem maiores que a chave, eles são deslocados uma posição para
a direita. Ao final, a chave é inserida na posição correta.

\begin{lstlisting}[caption={Implementação do Insertion Sort}]
void insertionSort(int* vetor, int tamanho) {
    for (int i = 1; i < tamanho; i++) {

        int chave = vetor[i];
        int j = i - 1;

        while (j >= 0 && vetor[j] > chave) {
            vetor[j + 1] = vetor[j];
            j--;
        }

        vetor[j + 1] = chave;
    }
}
\end{lstlisting}